\documentclass{claudeeditorial}

% Optional bibliography layer. Keep it in the document entrypoint, not in the
% class, so the visual system remains reusable for documents without citations.
\usepackage{csquotes}
\usepackage[backend=biber,style=authoryear,maxcitenames=2]{biblatex}
\addbibresource{bibliography/references.bib}

\renewcommand{\DocKicker}{Project Engineering}
\renewcommand{\DocTitle}{Manual de marca y uso}
\renewcommand{\DocSubtitle}{Documento diario para decisiones, riesgos, arquitectura, evidencias y seguimiento de ingeniería}
\renewcommand{\DocAuthor}{Equipo de proyecto}
\renewcommand{\DocRole}{Project engineering}
\renewcommand{\DocDate}{\today}
\renewcommand{\DocRef}{MAIN-001}
\renewcommand{\DocFooter}{claudeeditorial main}

\renewcommand{\ProjectName}{Nombre del proyecto}
\renewcommand{\ProjectPhase}{Fase actual}
\renewcommand{\ProjectOwner}{Equipo responsable}
\renewcommand{\ProjectStatus}{Activo}

\hypersetup{
  pdftitle={Manual de marca y uso},
  pdfauthor={Equipo de proyecto},
  pdfsubject={Diario de ingeniería de proyectos},
  pdfkeywords={project engineering, LaTeX, documentation, decision record}
}

\ClaudeRunningHeader
\setstretch{1.22}

\begin{document}
\BodyCopy

\ClaudeCover

\tableofcontents
\clearpage

\ClaudeChapterPage{01}{Situación diaria}{La página que mantiene alineado al equipo: qué cambió, qué se decidió, qué sigue abierto y dónde mirar.}

\ClaudeProjectHeader
  {\ProjectName}
  {Actualización ejecutiva y técnica para trabajo diario de ingeniería.}
  {\DocDate}
  {\ProjectStatus}

\section{Resumen ejecutivo}

\begin{claudehero}
\ClaudeLead{Escribe aquí la lectura de la jornada: el estado general, el movimiento más importante, la decisión que cambia el plan y la evidencia disponible. Mantén el texto breve y verificable.}
\end{claudehero}

\section{Indicadores}

\begin{tcbraster}[raster columns=4,raster equal height=rows,raster column skip=8pt]
\KPICard{Salud}{Verde}{sin bloqueo crítico}
\KPICard{Avance}{72\,\%}{incremento semanal}
\KPICard{Riesgos}{2}{uno requiere dueño}
\KPICard{Decisiones}{1}{ADR pendiente de firma}
\end{tcbraster}

\section{Decisión activa}

\ClaudeDecisionRecord[Pendiente]
  {ADR-001 · Título de la decisión}
  {Describe el contexto que obliga a decidir.}
  {Escribe la decisión concreta, en una frase que pueda auditarse.}
  {Explica el coste, la consecuencia y el siguiente control.}

\section{Riesgos}

\begin{ClaudeTable}{@{}lL L l@{}}
\ClaudeTableHeader Riesgo & Impacto & Mitigación & Estado \\
\toprule
\ClaudeRiskRow{R-001}{Riesgo técnico pendiente de validar.}{Añadir prueba reproducible y dueño.}{\StatusOpen}
\ClaudeRiskRow{R-002}{Dependencia externa con fecha incierta.}{Definir alternativa manual temporal.}{\StatusWIP}
\bottomrule
\end{ClaudeTable}

\section{Evidencia}

\begin{ClaudeCode}[title={Comprobación reproducible}]{SQL}
SELECT status, COUNT(*) AS total
FROM project_events
WHERE created_at >= CURRENT_DATE - 1
GROUP BY status
ORDER BY total DESC;
\end{ClaudeCode}

\begin{ClaudeOutput}
Registra aquí la salida esperada, una captura textual o la conclusión verificable.
\end{ClaudeOutput}

\section{Referencias}

El sistema de diseño sigue una disciplina de tokens cálidos y acento controlado, inspirada en reglas editoriales de productos de IA contemporáneos \parencite{claude-editorial-design-notes}.

\printbibliography

\ClaudeSignature{\DocAuthor}{\DocRole}

\end{document}
